\begin{table}[!ht]
\centering
\begin{threeparttable}
\caption{Summary Statistics}
\label{tab:sumstats}
\small
\setlength{\tabcolsep}{8pt}

\begin{tabularx}{\textwidth}{@{} l r r @{} }
\toprule
\textbf{Variable} & \textbf{Mean} & \textbf{SD} \\
\midrule
\multicolumn{3}{@{}l}{\textbf{Panel A: Outcome variables}} \\[2pt]
RN HPRD & 0.427 & 0.309 \\
LPN HPRD & 0.796 & 0.325 \\
CNA HPRD & 2.100 & 0.548 \\
Total HPRD & 3.324 & 0.793 \\
\addlinespace[0.8em]
\multicolumn{3}{@{}l}{\textbf{Panel B: Control variables}} \\[2pt]
Government (dummy) & 0.101 & 0.301 \\
Non-profit (dummy) & 0.171 & 0.377 \\
Chain affiliation (dummy) & 0.546 & 0.498 \\
Beds & 108.4 & 57.9 \\
Occupancy rate (\%) & 76.4 & 15.5 \\
\% Medicare & 13.0 & 11.1 \\
\% Medicaid & 60.5 & 20.2 \\
Acuity quartile 2 (state-month) & 0.231 & 0.422 \\
Acuity quartile 3 (state-month) & 0.230 & 0.421 \\
Acuity quartile 4 (state-month) & 0.233 & 0.423 \\
\bottomrule
\end{tabularx}

\begin{tablenotes}[flushleft]
\footnotesize
\item \textit{Notes:} The unit of observation is facility--month. HPRD denotes hours per resident day. RN, LPN, and CNA denote registered nurses, licensed practical nurses, and certified nursing assistants, respectively. Total HPRD is the sum of RN, LPN, and CNA HPRD. Occupancy rate is the ratio of residents to certified beds (percent). \% Medicare and \% Medicaid are the shares of residents covered by Medicare and Medicaid, respectively. Government and Non-profit are ownership-type indicator variables (for-profit omitted category). Chain affiliation is an indicator for membership in a multi-facility chain. Acuity quartiles are state--month quartiles of resident acuity (case-mix) with quartile 1 omitted. Rows $=$ 632,231; Facilities $=$ 7,806; Treated facilities $=$ 1,589; Period $=$ 2017/01–2024/06; Average months per facility $=$ 81.0.
\end{tablenotes}

\end{threeparttable}
\end{table}
